% Generated draft: requires mathematical and source review.
\documentclass[11pt]{article}
\usepackage[margin=25mm]{geometry}
\usepackage{fontspec}
\setmainfont{Latin Modern Roman}
\usepackage{amsmath,amssymb,mathtools,unicode-math}
\setmathfont{Latin Modern Math}
\usepackage{enumitem,needspace,xurl,hyperref}
\setlength{\parindent}{0pt}
\setlength{\parskip}{5pt}
\setlength{\emergencystretch}{3em}
\newcommand{\N}{\mathbb{N}}
\newcommand{\Z}{\mathbb{Z}}
\newcommand{\Q}{\mathbb{Q}}
\newcommand{\R}{\mathbb{R}}
\newcommand{\C}{\mathbb{C}}
\newcommand{\F}{\mathbb{F}}
\newcommand{\A}{\mathbb{A}}
\newcommand{\PP}{\mathbb{P}}
\newcommand{\E}{\mathbb{E}}
\newcommand{\eps}{\varepsilon}
\newcommand{\dd}{\,\mathrm{d}}
\begin{document}
\section*{UnsolvedMath ID: 701474}
\textbf{Source problem code:} catalogue entry
\subsection{Definitions and mathematical statement}
Let \(p\) be a prime and let \(G\) be a finite \(p\)-group. Write \([H,H]=\langle x^{-1}y^{-1}xy:x,y\in H\rangle\) for the derived subgroup of a subgroup \(H\leq G\). Define \(G^{(0)}=G\), \(G^{(i+1)}=[G^{(i)},G^{(i)}]\), and the derived length by
\[
\operatorname{dl}(G)=\min\{d\geq 0:G^{(d)}=1\}.
\]
The lower central series is defined by \(\gamma_1(G)=G\) and \(\gamma_{i+1}(G)=[\gamma_i(G),G]\). For nontrivial \(G\), its nilpotency class is the least \(c\) such that \(\gamma_{c+1}(G)=1\). The maximal-class hypothesis means that \(|G|=p^n\) and its nilpotency class is \(n-1\).

Interpret irreducible characters in the usual ordinary complex sense. Assume that
\[
\forall\chi\in\operatorname{Irr}_{\C}(G)\quad
\exists N\trianglelefteq G\quad
\exists\lambda\in\operatorname{Irr}_{\C}(N):\qquad
\lambda(1)=1\quad\text{and}\quad\chi=\operatorname{Ind}_{N}^{G}\lambda.
\]
Here \(\operatorname{Irr}_{\C}(H)\) denotes the set of irreducible complex characters of \(H\), and \(\operatorname{Ind}_{N}^{G}\) denotes character induction.

Let \(\mathcal C(G)\) denote the collection of subgroups occurring in the central series intended by the source, and put
\[
S_{\mathcal C}(G)=\{[H,H]:H\in\mathcal C(G)\}.
\]
This is a set of distinct subgroups, so repetitions are counted only once. The source does not identify which central series is intended. In particular, it does not specify the lower central series or the upper central series, the latter being defined by \(Z_0(G)=1\) and \(Z_{i+1}(G)/Z_i(G)=Z(G/Z_i(G))\), where \(Z(H)\) is the center of \(H\).

The question asks whether the derived length has a logarithmic bound in terms of \(|S_{\mathcal C}(G)|\). Once the intended series is fixed, a possible precise interpretation is the existence of constants \(A>0\) and \(B\geq0\) such that
\[
\operatorname{dl}(G)\leq A\log\bigl(1+|S_{\mathcal C}(G)|\bigr)+B
\]
for every group satisfying the hypotheses, with \(\log\) the natural logarithm. A fully uniform interpretation requires the same constants for every prime \(p\) and every such \(G\); a fixed-prime interpretation permits constants depending on \(p\). The source does not resolve this distinction or specify the constants or exact form of the logarithmic bound. These qualifications remain unresolved in this statement-only draft.
\subsection{Short English statement}
T. M. Keller asks whether a finite \(p\)-group of maximal class, whose irreducible complex characters all arise by inducing linear characters from normal subgroups, has derived length bounded logarithmically by the number of distinct derived subgroups of terms in its central series. The intended central series and the allowed dependence of the bound on \(p\) require clarification.
\subsection{Sources}
\begin{itemize}
\item Supplied local UnsolvedMath registry record, permanent ID \(701474\), exact source problem code \texttt{catalogue entry}, title \emph{Kourovka Notebook Problem 18.45}; upstream revision \texttt{d371ee917ac9647fbb06cd6fc529deec6f3dda82}. The record was created and updated on 2026-10-16 and labels the problem \emph{open}; this is a supplied registry label, not an independently verified mathematical or literature assessment.
\item T. M. Keller, Problem 18.45, \emph{The Kourovka Notebook (Unsolved Problems in Group Theory)}, 18th issue (2014), as attributed in the supplied record. The background also identifies the 21st issue edition (2026), \texttt{arXiv:1401.0300v48}. The supplied collection description names E. I. Khukhro and V. D. Mazurov as editors. These bibliographic details are supplied attribution and have not been independently checked.
\end{itemize}
\end{document}
